\documentclass[leqno]{jsarticle}
\usepackage{amssymb}
\usepackage{amsthm}
\usepackage{amsmath}
\usepackage{comment}
%定理環境 thmはあまり使わずpropをなるべく使う。
%主張は基本的に証明の中で使い、そのため番号を付けない。
%注意環境を付けてみた。
\newtheorem{thm}{定理}
\newtheorem{prop}[thm]{命題}
\newtheorem{cor}[thm]{系}
\newtheorem{lem}[thm]{補題}
\newtheorem{df}[thm]{定義}
\newtheorem{exam}[thm]{例}
\newtheorem{fact}[thm]{事実}
\newtheorem*{claim}{主張}
\newtheorem*{cau}{注意}
\renewcommand{\proofname}{{\bf 証明}}
%矢印たち。
%\toは\rightarrowと同じ。\getsは\leftarrowと同じ。同値は\iff
\newcommand{\To}{\Longrightarrow}
\newcommand{\Gets}{\LongLeftarrow}
%黒板太字
\newcommand{\nn}{\mathbb{N}}
\newcommand{\zz}{\mathbb{Z}}
\newcommand{\qq}{\mathbb{Q}}
\newcommand{\rr}{\mathbb{R}}
\newcommand{\cc}{\mathbb{C}}
%無限大
\newcommand{\mgn}{\infty}
%差集合は\setminusで統一する。等号の否定は\neq
%真部分集合は\subsetneqq　偏微分は\partial
%ディスプレイスタイル。\dfracでディスプレイ分数
\newcommand{\dpst}{\displaystyle}

\newcommand{\cho}[2]{#1 \choose #2}
\newcommand{\nik}[2]{{}_{#1}\text{C}_{#2}}
\newcommand{\Be}[3]{B_{#1,#2}(#3)}
\newcommand{\bb}{B_{n,k}(x)}
\newcommand{\sss}{D_{n,k}(x,y)}
\newcommand{\dd}{\partial}
\title{多項式近似定理}
\author{電波通信}
\date{\today}
\begin{document}
\maketitle
この文章では実直線の有界閉区間上の連続関数が多項式で一様に近似出来るというワイエルストラスの多項式近似定理をベルンシュタイン多項式を用いた方法で証明する。多項式近似定理は一般の有界閉区間上の関数に対する定理だが、一般の有界閉区間は$[0,1]$と同相なので、$[0,1]$上で証明してしまっても構わない。\\ 
早速証明に必要な補題を準備しよう。多項式$\bb$を
\[
B_{n,k}(x):=\binom nk x^k(1-x)^{n-k}=\nik{n}{k}\ x^k(1-x)^{n-k}\ \ (0\le k\le n)
\]
と定義する。$\binom nk =\nik{n}{k}$は二項係数である。$[0,1]上この多項式は常に0以上となることに注意せよ。$
このとき次が成立する。

\begin{lem}${}$
\begin{eqnarray}
\displaystyle &\sum_{k=0}^{n}\bb=1\\ 
\displaystyle &\sum_{k=0}^{n}k\bb=nx\\ 
\displaystyle &\sum_{k=0}^{n}k(k-1)\bb=n(n-1)x\\ 
\displaystyle &\sum_{k=0}^{n}k^2\bb=n(n-1)x^2+nx\\ 
\displaystyle &\sum_{k=0}^{n}(nx-k)^2\bb=nx(1-x)
\end{eqnarray}
\end{lem}
\begin{proof}
$\bb$を２変数化して証明する。
\[
\sss := \binom nk x^ky^{n-k}=\nik{n}{k}\ x^ky^{n-k}
\]とおけば、$\bb=D_{n,k}(x,1-x)$となる。まず1番の式に関しては二項係数の定義から
\[
\sum_{k=1}^{n}\sss=(x+y)^n
\]から$y=1-x$を代入すればよい。\\ 
2番の式に関しては微分公式$\frac{d}{dx}x^k=kx^{k-1}$から$kx^k=x\frac{d}{dx}x^k$なので
\[
\sum_{k=1}^{n}k\sss=x\frac{\dd}{\dd x}\sum_{k=1}^{n}\sss=x\frac{\dd}{\dd x}(x+y)^n=nx(x+y)^{n-1}
\]となり、やはり$y=1-x$を代入すればわかる。\\ 
第3式も同様に$k(k-1)x^{k}=x^2\frac{d^2}{dx^2}x^k$から
\[
\sum_{k=1}^{n}k(k-1)\sss=x^2\frac{\dd^2}{\dd x^2}\sum_{k=1}^{n}\sss=x^2\frac{\dd^2}{\dd x^2}(x+y)^n=n(n-1)x^2(x+y)^{n-2}
\]なのでやっぱり$y=1-x$を代入すれば良い。\\ 
第4,5式に関しては展開して1,2,3式を用いれば容易にわかる。
\end{proof}

さていよいよ多項式近似定理を証明しよう。以下では$||\bullet||_{\infty}$は$[0,1]$上の一様ノルム（$\sup$ノルム）である。
\begin{thm}[多項式近似定理]
$[0,1]$上の連続関数$f$に対して
\[p_n(x)=\sum_{k=0}^{n}f\left(\frac{k}{n}\right)\Be{n}{k}{x}\]と定義すると
$||f-p_n||_{\infty}\rightarrow 0$ as $ n\rightarrow \infty$が成立する。
\end{thm}この多項式$p_n$を$f$のベルンシュタイン多項式と言う。\\ 
\proofname\\ 
\[
\forall \varepsilon >0\ \exists N\in \mathbb{N}\ n>N\Rightarrow ||f-p_n||_{\infty}\le\varepsilon
\]を示す。
コンパクト集合上の連続関数は一様連続であるから
\[
\forall \varepsilon>0\ \exists \delta>0\ \forall x,y\in [0,1]\ |x-y|<\delta \Rightarrow |f(x)-f(y)|<\frac{\varepsilon}{2}
\]
となる。$\displaystyle M=\max _{x\in M}f(x)$とし、さきの$\delta$に対して$N$を
\[
n>N\Rightarrow \frac{M}{2\delta^2}\frac{1}{n}<\frac{\varepsilon}{2}
\]となるように取る。さて、任意の$x\in [0,1]$に対して
\[
|f(x)-p_n(x)|\le\sum_{|x-\frac{k}{n}|<\delta}\left|f(x)-f\left(\frac{k}{n}\right)\!\right|\bb+\sum_{|x-\frac{k}{n}|\ge\delta}\left|f(x)-f\left(\frac{k}{n}\right)\!\right|\bb
\]
となる。ここで$\displaystyle \sum_{|x-\frac{k}{n}|<\delta}$の意味は$|x-\frac{k}{n}|$を充たす$k$について和を取ると言う意味である。(もちろん$0\le k\le n$である。)ただし$|x-\frac{k}{n}|$となる$k$が存在しなければ$0$と定義する。$\displaystyle \sum_{|x-\frac{k}{n}|\ge\delta}$も同様の意味である。
まず上式の初項について一様連続性から$\displaystyle{\left|f(x)-f\left(\frac{k}{n}\right)\!\right|<\frac{\varepsilon}{2}}$となり
\[
\sum_{|x-\frac{k}{n}|<\delta}\left|f(x)-f\left(\frac{k}{n}\right)\!\right|\bb<\frac{\varepsilon}{2}\sum_{|x-\frac{k}{n}|<\delta}\bb\le \frac{\varepsilon}{2}\sum_{k=1}^{n}\bb=\frac{\varepsilon}{2}
\]
ここで補題の１番の式を使った。\\ 
さて第二項について$|x-\frac{k}{n}|\ge\delta$より$|nx-k|\ge n\delta$よって$(nx-k)^2\ge n^2\delta^2$これを元に
\[
n^2\delta^2\sum_{|x-\frac{k}{n}|\ge\delta}\left|f(x)-f\left(\frac{k}{n}\right)\!\right|\bb\le 2M\sum_{|x-\frac{k}{n}|\ge\delta}(nx-k)^2\bb\le 2M\sum_{k=1}^{n}(nx-k)^2\bb\]
$
\displaystyle{=2Mnx(1-x)\le 2Mn\frac{1}{4}=M\frac{n}{2}}
$を得る。ここで$|f(x)-f(y)|\le 2M$と補題の第５式を用いた。\\ 
よって$n>N$ならば
\[
\sum_{|x-\frac{k}{n}|\ge\delta}\left|f(x)-f\left(\frac{k}{n}\right)\!\right|\bb\le \frac{M}{2\delta^2}\frac{1}{n}
<\frac{\varepsilon}{2}
\]
以上より$n>N$ならば
\[
|f(x)-p_n(x)|<\frac{\varepsilon}{2}+\frac{\varepsilon}{2}=\varepsilon
\]
この$N$は$x$に依存せず、$x$は任意であったから$n>N\Rightarrow ||f-p_n||_{\infty}\le\varepsilon$つまり$||f-p_n||_{\infty}\rightarrow 0$ as $ n\rightarrow \infty$ $\square$



\begin{thebibliography}{9}
\bibitem{1}杉浦光夫『解析入門I』（基礎数学２）東京大学出版会(1980)
\end{thebibliography}






















\end{document}